\renewcommand{\abstractname}{Resumen}
\begin{abstract}
Este trabajo investiga el potencial de la computación cuántica para la optimización de hiperparámetros (HPO) en modelos de aprendizaje profundo, centrándose específicamente en arquitecturas de tipo Transformer aplicadas al monitoreo predictivo de procesos (PPM) con el dataset real \textit{env\_permit}. El problema se modela formalmente como una Optimización Binaria Cuadrática sin Restricciones (QUBO) de exactamente 30 variables (30 cúbits) con restricciones One-Hot, mapeando un espacio de búsqueda de $2^{30}$ combinaciones de hiperparámetros. El objetivo principal es explorar la viabilidad de resolvedores cuánticos variacionales bajo el ruido físico de la era NISQ, comparando la Optimización Bayesiana clásica frente al algoritmo resolvedor cuántico variacional (VQE) estándar y la evolución en tiempo imaginario cuántico variacional (VarQITE).

\vspace{0.5cm}
Mediante simulaciones numéricas en el supercomputador FinisTerrae III y su posterior validación en la unidad de procesamiento cuántico (QPU) superconductora real Qmio (32 cúbits) del CESGA, se evaluó experimentalmente el comportamiento físico de los algoritmos ante la decoherencia, errores de lectura y ruido térmico. Los resultados revelan una paradoja física y temporal crucial: mientras VQE (sintonizado secuencialmente con COBYLA) fracasa en la QPU real con una tasa de factibilidad del 0\% debido al impacto de mínimos locales y ruido, VarQITE alcanza una factibilidad del 100\% sin violar ninguna restricción. Esta resiliencia se debe al efecto de filtrado exponencial ($e^{-E_n\tau}$) en tiempo imaginario que extingue las componentes no factibles mediante la penalización de Lagrange ($\lambda$). Adicionalmente, VarQITE es un orden de magnitud más rápido en hardware real (~220s frente a ~1500s de VQE) gracias a que su esquema de evolución por pasos fijos permite el envío de circuitos en paralelo (\textit{batching}) minimizando las peticiones de red física al criostato, invirtiendo la tendencia de la simulación clásica clásica. El estudio valida empíricamente la efectividad de VarQITE en la frontera de Pareto de HPO, demostrando el potencial de la computación cuántica real para el diseño de Inteligencia Artificial de alta escala.
\end{abstract}
